\documentclass[a4paper,12pt]{article}
\usepackage[T2A]{fontenc}
\usepackage[utf8]{inputenc}
\usepackage[russian]{babel}
\usepackage{amsmath,amssymb,mathtools}
\usepackage{icomma}
\usepackage[a4paper,left=23mm,right=23mm,top=22mm,bottom=23mm]{geometry}
\usepackage{microtype}
\usepackage{tikz}
\usetikzlibrary{calc,arrows.meta,positioning,shapes.geometric}
\usepackage{tabularx,booktabs}
\usepackage{enumitem}
\usepackage{xcolor}
\usepackage{fancyhdr}
\usepackage[hidelinks]{hyperref}
\definecolor{accent}{HTML}{087D80}
\definecolor{darkaccent}{HTML}{006166}
\definecolor{softgray}{HTML}{585D62}
\definecolor{textgray}{HTML}{24272A}
\hypersetup{pdftitle={Чек-лист по физике. Наклонная тяга, трение и второй закон Ньютона},pdfauthor={Лёвин Артём Александрович}}
\setlength{\parindent}{0pt}
\setlength{\parskip}{5.5pt}
\setlength{\headheight}{25pt}
\renewcommand{\arraystretch}{1.25}
\setlist[enumerate]{leftmargin=*,itemsep=.45em,topsep=.25em}
\setlist[itemize]{leftmargin=*,itemsep=.3em,topsep=.25em}
\color{textgray}
\pagestyle{fancy}
\fancyhf{}
\fancyhead[L]{\small\color{softgray}Наклонная тяга и трение}
\fancyhead[R]{\small\color{softgray}Володя Хачатурян}
\fancyfoot[C]{\small\thepage}
\renewcommand{\headrulewidth}{.5pt}
\renewcommand{\headrule}{\hbox to\headwidth{\color{accent}\leaders\hrule height \headrulewidth\hfill}}
\newcommand{\sect}[1]{\section{#1}\vspace{-1.0em}\noindent\textcolor{accent}{\rule{\linewidth}{.5pt}}\par\vspace{.4em}}
\newcommand{\subsect}[1]{\vspace{.35em}{\color{darkaccent}\bfseries #1}\par\vspace{-.1em}}
\newcommand{\key}[1]{{\color{darkaccent}\bfseries #1}}
\newcommand{\unit}[1]{\,\text{#1}}
\begin{document}
\pagenumbering{gobble}
\thispagestyle{empty}
\begin{center}
\vspace*{24mm}
{\large\color{darkaccent} ФИЗИКА}\par
\vspace{13mm}
{\LARGE Чек-лист}\par
\vspace{13mm}
{\Large\color{accent} Силы при наклонной тяге:}\par
{\Large\color{accent} проекции, трение, ускорение}\par
\vspace{15mm}
{\normalsize Подготовка к ОГЭ по физике}\par
{\normalsize Повторение по материалу занятия 4 ноября 2026 г.}\par
\vfill
{\large Володя Хачатурян}\par
\vspace{9mm}
{9 ноября 2026 г.}\par
\vspace{14mm}
{\small Лёвин Артём Александрович}\par
{\small эксклюзивно для Володи Хачатуряна}\par
\vspace{17mm}
\begin{tikzpicture}[x=1cm,y=1cm]
\draw[accent,line width=.8pt] (0,0) .. controls (1,.40) and (2,-.39) .. (3,0)
 .. controls (4,.4) and (5,-.38) .. (6,0)
 .. controls (7,.35) and (8,-.38) .. (9,0);
\end{tikzpicture}
\end{center}
\clearpage
\pagenumbering{arabic}
\setcounter{page}{1}

\sect{Как увидеть физику в условии}

\key{Главная идея.} При тяге под углом $30^\circ$ к горизонтали сила $F$ выполняет две роли: её горизонтальная составляющая тянет брусок, а вертикальная приподнимает его и уменьшает силу реакции опоры $N$. Из-за этого меняется и сила трения.

Рассматриваем \textbf{брусок, скользящий вправо} по горизонтальному столу. Сила тяги направлена вправо-вверх, ускорения по вертикали нет. Сопротивлением воздуха пренебрегаем.

\begin{center}
\begin{tikzpicture}[scale=.94,>=Latex]
\draw[softgray,line width=.85pt] (-3,-.7)--(5,-.7);
\foreach \x in {-2.8,-2.35,...,4.8}{\draw[softgray!45,line width=.3pt] (\x,-.7)--++(-.13,-.14);}
\fill[accent!7] (-.75,-.7) rectangle (.75,.62);
\draw[softgray,line width=.75pt] (-.75,-.7) rectangle (.75,.62);
\node at (0,-.13) {$m$};
\coordinate (O) at (0,-.1);
\draw[->,line width=1.0pt] (O)--(0,1.7) node[above] {$\vec N$};
\draw[->,line width=1.0pt] (O)--(0,-1.85) node[below] {$m\vec g$};
\draw[->,line width=1.0pt] (O)--(-2.4,-.1) node[left] {$\vec F_{\text{тр}}$};
\draw[->,accent,line width=1.45pt] (O)--(2.65,1.43) node[above right] {$\vec F$};
\draw[softgray,dashed] (O)--(2.65,-.1) node[pos=.72,below] {$F_x$};
\draw[softgray,dashed] (2.65,-.1)--(2.65,1.43) node[pos=.55,right] {$F_y$};
\draw[accent] (.77,-.1) arc[start angle=0,end angle=30,radius=.77];
\node[accent] at (1.26,.19) {$30^\circ$};
\draw[->,softgray] (4,-1.55)--(4.8,-1.55) node[right] {$x$};
\draw[->,softgray] (4,-1.55)--(4,-.77) node[above] {$y$};
\end{tikzpicture}
\end{center}
\vspace{-2mm}
\subsect{На брусок действуют четыре силы}
\begin{itemize}
\item $m\vec g$ --- сила тяжести, направлена вниз;
\item $\vec N$ --- реакция опоры, направлена вверх;
\item $\vec F$ --- наклонная сила тяги, направлена вправо-вверх;
\item $\vec F_{\text{тр}}$ --- сила трения скольжения, направлена влево, против скольжения.
\end{itemize}

\sect{Разложение силы при угле 30 градусов}

Вертикальная проекция лежит \textbf{против угла $30^\circ$} в прямоугольном треугольнике и равна половине гипотенузы. Горизонтальную проекцию найдём по теореме Пифагора:
\[
\boxed{F_y=\frac{F}{2}},\qquad
F_x=\sqrt{F^2-F_y^2}=\boxed{\frac{F\sqrt3}{2}}.
\]

\subsect{Пример}
При $F=18\unit{Н}$:
\[
F_y=9\unit{Н},\qquad
F_x=\sqrt{18^2-9^2}=9\sqrt3\unit{Н}.
\]
\key{Проверка:} $F_x^2+F_y^2=F^2$; обе проекции меньше модуля $F$.

\clearpage
\sect{Второй закон Ньютона по осям}

\subsect{По вертикали}
Брусок не ускоряется вверх или вниз, поэтому
\[
N+F_y-mg=0 \quad\Longrightarrow\quad
\boxed{N=mg-\frac{F}{2}}.
\]

\subsect{Трение скольжения}
В школьной модели его модуль равен
\[
\boxed{F_{\text{тр}}=\mu N=\mu\left(mg-\frac{F}{2}\right)}.
\]
Здесь $\mu$ --- коэффициент трения скольжения, без единиц измерения.

\subsect{По горизонтали}
Положительное направление оси $x$ --- вправо, поэтому
\[
\boxed{F_x-F_{\text{тр}}=ma_x}.
\]
Объединяя формулы, получаем полезное выражение:
\[
\boxed{a_x=\frac{\dfrac{F\sqrt3}{2}-\mu\left(mg-\dfrac F2\right)}{m}}.
\]
\key{Граница применимости:} брусок должен соприкасаться со столом: $N>0$, то есть $F<2mg$. Если расчёт дал $N\le0$, привычная схема скольжения по столу требует пересмотра.

\subsect{Что означает знак ускорения}
При скольжении вправо $a_x>0$ означает разгон; $a_x<0$ означает замедление. После остановки направление трения и режим движения необходимо определять заново.

\sect{Три приёма решения задач}

\subsect{Приём 1. Найти ускорение}
Брусок массой $3$ кг скользит вправо. Сила тяги $18$ Н направлена под углом $30^\circ$, $\mu=0{,}2$; примем $g=10\,\text{м/с}^2$.
\[
N=3\cdot10-\frac{18}{2}=21\unit{Н},\qquad
F_{\text{тр}}=0{,}2\cdot21=4{,}2\unit{Н};
\]
\[
F_x=9\sqrt3\unit{Н},\qquad
a_x=\frac{9\sqrt3-4{,}2}{3}
=3\sqrt3-1{,}4\approx\boxed{3{,}80\,\text{м/с}^2}.
\]

\subsect{Приём 2. Найти массу по известному трению}
На брусок действует тяга $F=16$ Н; $\mu=0{,}25$, трение скольжения равно $8$ Н. Из $F_{\text{тр}}=\mu N$:
\[
N=\frac{8}{0{,}25}=32\unit{Н},\quad
mg=N+F_y=32+8=40\unit{Н},\quad
\boxed{m=4\unit{кг}}.
\]

\subsect{Приём 3. Найти неизвестную силу тяги}
Из уравнений по двум осям при известном ускорении $a_x$:
\[
\frac{F\sqrt3}{2}-\mu\left(mg-\frac F2\right)=ma_x,
\]
\[
\boxed{F=\frac{2m(a_x+\mu g)}{\sqrt3+\mu}}.
\]
После вычисления $F$ обязательно проверяем условие контакта $F<2mg$.


\sect{Типичные ошибки и итоговый самоконтроль}
\begin{enumerate}
\item \key{Вдоль оси $x$ подставили всю силу $F$.} Используйте горизонтальную проекцию $F_x$, поскольку тяга наклонена.
\item \key{Записали $N=mg$.} Вертикальная составляющая тяги разгружает опору: $N=mg-F/2$.
\item \key{Перепутали катеты.} При угле $30^\circ$ от горизонтали $F_y=F/2$, а $F_x=F\sqrt3/2$.
\item \key{Забыли знак трения.} При движении вправо проекция трения отрицательна.
\item \key{Исправили отрицательное ускорение на положительное.} Отрицательный знак показывает направление, а при движении вправо --- замедление.
\item \key{Получили отрицательную реакцию опоры.} Прежде чем применять $\mu N$, проверьте возможность контакта бруска со столом.
\end{enumerate}

\subsect{Проверяю себя перед сдачей}
\begin{itemize}
\item[$\square$] Нарисованы и подписаны все четыре силы.
\item[$\square$] Угол отсчитан от горизонтали; проекции записаны правильно.
\item[$\square$] В вертикальном балансе стоит нуль, поскольку $a_y=0$.
\item[$\square$] Реакция опоры положительна: $N>0$.
\item[$\square$] Сила трения найдена через $\mu N$.
\item[$\square$] При записи второго закона Ньютона учтены знаки проекций.
\item[$\square$] Ответ содержит единицы измерения и проверен по смыслу.
\end{itemize}

\clearpage
\begin{center}
{\LARGE\color{darkaccent} Алгоритм решения задачи}\par
\vspace{2mm}
{\small Брусок скользит вправо; сила тяги направлена под углом $30^\circ$ к горизонтали}\par
\vspace{12mm}
\begin{tikzpicture}[
 >=Latex,
 node distance=9mm,
 beginend/.style={ellipse,draw=darkaccent,line width=.85pt,align=center,text width=54mm,minimum height=12mm,font=\small},
 flowstep/.style={rectangle,rounded corners=2mm,draw=accent,line width=.8pt,align=center,text width=118mm,minimum height=18mm,inner sep=2mm,font=\small},
 decision/.style={diamond,aspect=3.2,draw=accent,line width=.9pt,align=center,text width=42mm,inner sep=1mm,font=\small},
 arr/.style={-{Latex[length=2.5mm]},draw=darkaccent,line width=.95pt}]
\node[beginend] (s) {Начало};
\node[flowstep,below=of s] (forces) {Схема сил: $m\vec g$, $\vec N$, $\vec F$, $\vec F_{\text{тр}}$; оси $x$ вправо, $y$ вверх};
\node[flowstep,below=of forces] (project) {Проекции тяги: $F_y=F/2$, $F_x=F\sqrt3/2$};
\node[flowstep,below=of project] (normal) {Вертикальный баланс: $N=mg-F_y$};
\node[decision,below=12mm of normal] (contact) {$N>0$?};
\node[beginend,text width=43mm,minimum height=16mm,right=5mm of contact] (stop) {Конец: модель трения по столу неприменима};
\node[flowstep,below=14mm of contact] (fric) {Трение: $F_{\text{тр}}=\mu N$; горизонтальный баланс: $F_x-F_{\text{тр}}=ma_x$};
\node[flowstep,below=of fric] (ans) {Выразить искомое, вычислить, проверить единицы и знак результата};
\node[beginend,below=of ans] (end) {Ответ};
\draw[arr] (s)--(forces);
\draw[arr] (forces)--(project);
\draw[arr] (project)--(normal);
\draw[arr] (normal)--(contact);
\draw[arr] (contact)--node[right,font=\small]{да}(fric);
\draw[arr] (contact.east)--node[above,font=\small]{нет}(stop.west);
\draw[arr] (fric)--(ans);
\draw[arr] (ans)--(end);
\end{tikzpicture}
\end{center}
\clearpage

\clearpage
\section*{Тренировочный блок. Путь А}
\noindent\textcolor{accent}{\rule{\linewidth}{.5pt}}\par
Во всех задачах $g=10\,\text{м/с}^2$, угол тяги к горизонтали равен $30^\circ$, тяга направлена вправо-вверх. Если сказано о скольжении, брусок скользит вправо и трение действует влево. Предполагайте наличие контакта с опорой.

\subsect{Базовый уровень}
\begin{enumerate}[label=\textbf{\arabic*.}]
\item На брусок действует наклонная сила $F=16$ Н. Найдите её проекции $F_x$ и $F_y$.
\item Масса бруска $3$ кг, модуль наклонной тяги $10$ Н. Найдите реакцию опоры.
\item Брусок массой $2$ кг скользит вправо. Сила тяги $8$ Н, $\mu=0{,}25$. Найдите модуль силы трения скольжения.
\end{enumerate}

\subsect{Средний уровень}
\begin{enumerate}[label=\textbf{\arabic*.},start=4]
\item Брусок массой $4$ кг скользит вправо. Сила тяги $12$ Н, $\mu=0{,}3$. Найдите $N$ и $F_{\text{тр}}$.
\item Модуль тяги $18$ Н, $\mu=0{,}2$, сила трения скольжения $4{,}2$ Н. Найдите массу бруска.
\item Брусок массой $5$ кг скользит вправо. Сила тяги $20$ Н, $\mu=0{,}25$. Найдите проекцию ускорения $a_x$.
\item Брусок массой $3$ кг скользит вправо. Сила тяги $12$ Н, $\mu=0{,}4$. Найдите проекцию ускорения $a_x$.
\end{enumerate}

\subsect{Повышенный уровень}
\begin{enumerate}[label=\textbf{\arabic*.},start=8]
\item Брусок массой $4$ кг скользит вправо. Сила тяги $14$ Н, трение скольжения равно $6{,}6$ Н. Найдите коэффициент трения $\mu$.
\item Брусок массой $2$ кг скользит вправо с ускорением $2\,\text{м/с}^2$. Коэффициент трения $\mu=0{,}2$. Найдите модуль силы тяги $F$.
\item Брусок массой $4$ кг скользит вправо. Сила тяги $10$ Н, $\mu=0{,}6$. Найдите $a_x$ и объясните физический смысл знака.
\end{enumerate}

\clearpage
\section*{Ответы к тренировочному блоку}
\noindent\textcolor{accent}{\rule{\linewidth}{.5pt}}\par
\vspace{3mm}
\renewcommand{\arraystretch}{1.6}
\begin{tabularx}{\linewidth}{@{}>{\bfseries}p{10mm}X@{}}
\toprule
№ & Ответ\\
\midrule
1 & $F_x=8\sqrt3$ Н, $F_y=8$ Н.\\
2 & $N=25$ Н.\\
3 & $F_{\text{тр}}=4$ Н.\\
4 & $N=34$ Н, $F_{\text{тр}}=10{,}2$ Н.\\
5 & $m=3$ кг.\\
6 & $a_x=2\sqrt3-2\approx1{,}46\,\text{м/с}^2$.\\
7 & $a_x=2\sqrt3-3{,}2\approx0{,}26\,\text{м/с}^2$.\\
8 & $\mu=0{,}2$.\\
9 & $F=\dfrac{16}{\sqrt3+0{,}2}\approx8{,}28$ Н.\\
10 & $a_x=\dfrac{5\sqrt3-21}{4}\approx-3{,}08\,\text{м/с}^2$; ускорение направлено влево, до остановки брусок замедляется.\\
\bottomrule
\end{tabularx}

\vspace{12mm}
\subsect{Итоговая формула для самопроверки}
\[
\boxed{\;F_x-F_{\text{тр}}=ma_x,\qquad N=mg-\frac F2,\qquad F_{\text{тр}}=\mu N.\;}
\]

\vfill
\begin{center}\small\color{softgray}
Лёвин Артём Александрович --- эксклюзивно для Володи Хачатуряна
\end{center}
\end{document}